% Conservative incorporation: additive selection and thermal instrument.
\section{Additive selection and the thermal reading instrument}
\label{sec:ii-aditiva-instrumento}

The spectral reference and action constructed in the preceding sections
make it possible to specify what is counted when information and energy
are measured. The thermal coefficient appears as a relative intensity
between two events in the same record. Additivity characterizes its
informational role, and both events admit an explicit apparatus preserving
the complete state. This develops the marked rule of
\ref{thm:ii-ter27-transducer}, including its realization, its selection
within a class of readers, and its transport conditions.

The composition retains its APP--TRIT--TPK antecedent: additive and
multiplicative sheets, regime, orientation, remainder, quotient, boundary,
route, and memory. The prolongation
$w_6\to w_{12}\to w_{18}\to w_{24}\to w_{30}\to R_{36}\to G_9$
and the order $U_t\to\Gamma_9\to K_{\rm ph}$ precede this spectral
reading. Phase return prolongs the record; the observed marks replace
neither that record nor the two projections of the joint discrete structure
of the continuum. The ratio $r=1/10$, the ranks $(1,380,649)$, and their
assignment to degrees $23+3n$ are received with the conditions of
\ref{subsec:ii-ter27-reference}. This assignment remains an explicit
part of the realization, together with the Hamiltonian and its energy origin.

% BEGIN THERMAL_R03_01
\subsection{Spectral events and relative measure}
\label{subsec:ii-tins-eventos}

Let $\mathcal K=\ell^2(\mathbb N_0)\otimes\mathbb C^M$, with $M=649$,
and consider the normalized density
\begin{equation}
 \Sigma=\rho_r\otimes I_M/M,\qquad
 \rho_r=(1-r)\sum_{j\ge0}r^j|j\rangle\langle j|.
 \label{eq:ii-tins-sigma}
\end{equation}
For $n=0,1,2$, let $j_n=23+3n$ and let $Q_n$ be memory projectors of rank
$d_n$, where $(d_0,d_1,d_2)=(1,380,649)$. Fix a unit vector $\chi_B$.
The letters $A,B,C$ here denote reading events, not geometric area:
\begin{align}
 P_A&=\sum_{n=0}^2|j_n\rangle\langle j_n|\otimes Q_n,&
 P_B&=|0\rangle\langle0|\otimes|\chi_B\rangle\langle \chi_B|,
 \nonumber\\
 P_C&=I-P_A-P_B,& \mu(P)&=\Tr(\Sigma P).
 \label{eq:ii-tins-marks}
\end{align}
The three projectors are mutually orthogonal, sum to the identity, and
commute with $\Sigma$. Event $B$ fixes a rank-one reference; $C$ retains
the entire complement of the marked sectors.

\Needspace{12\baselineskip}
\begin{samepage}
\begin{proposition}[Relative intensity and compression]
\label{prop:ii-tins-relative}
Writing $p_A=\mu(P_A)$ and $p_B=\mu(P_B)$ gives
\begin{equation}
 p_B=\frac{1-r}{M},\qquad
 p_A=\frac{1-r}{M}\sum_{n=0}^2d_nr^{23+3n},\qquad
 \frac{p_A}{p_B}=b_*=
 \frac{1380649}{10^{29}}=1.380649\times10^{-23}.
 \label{eq:ii-tins-relative}
\end{equation}
If $\iota:\mathcal D\to\mathcal K$ embeds each carrier sector in its
degree $j_n$ and in $\operatorname{ran}Q_n$, then
\begin{equation}
 p_B^{-1}\iota^*\Sigma\iota
 =\bigoplus_{n=0}^2r^{23+3n}I_{d_n}=\eta.
 \label{eq:ii-tins-compression}
\end{equation}
\end{proposition}
\end{samepage}
\begin{proof}
At degree $j$, the density acts as $(1-r)r^j/M$ times the memory
identity. Taking its trace over a subspace of rank $d$ multiplies that
scalar by $d$. Applying this to degree zero and the three marked degrees
gives the two probabilities; their ratio cancels the common factor.
Since $r=10^{-1}$,
$\sum_nd_nr^{23+3n}=10^{-23}(1+380\,10^{-3}+649\,10^{-6})$,
which yields the displayed fraction. The same scalar action, restricted
by $\iota$, proves \eqref{eq:ii-tins-compression}.
\end{proof}

The intensity $b_*$ is a relative measure, not the raw probability $p_A$.
Using the whole vacuum $|0\rangle\langle0|\otimes I_M$ as reference
would give $b_*/M$: this changes the reference event. By contrast, adjoining
an independent auxiliary factor and extending both projectors by its
identity preserves the two probabilities and their ratio.
% END THERMAL_R03_01

% BEGIN THERMAL_R03_02
\subsection{Unitary apparatus, recovery, and energy domain}
\label{subsec:ii-tins-aparato}

On a three-dimensional pointer with basis $|A\rangle,|B\rangle,|C\rangle$,
define
\begin{equation}
 V\psi=\sum_{y=A,B,C}P_y\psi\otimes|y\rangle.
 \label{eq:ii-tins-isometry}
\end{equation}
For a unitary realization, let $R_C=I$, let $R_A$ exchange $|C\rangle$
and $|A\rangle$, and let $R_B$ exchange $|C\rangle$ and $|B\rangle$.
The coupling is
\begin{equation}
 \mathbb U=\sum_{y=A,B,C}P_y\otimes R_y,
 \qquad \mathbb U(\psi\otimes|C\rangle)=V\psi.
 \label{eq:ii-tins-unitary}
\end{equation}

\begin{proposition}[Conservative recording of the three alternatives]
\label{prop:ii-tins-instrument}
The identities $\mathbb U^*\mathbb U=\mathbb U\mathbb U^*=I$ and
$V^*V=I$ hold. For any density $\xi$, the joint state
$\widehat\xi=V\xi V^*$ admits exact recovery $V^*\widehat\xi V=\xi$.
Reading the pointer defines the instrument
\begin{equation}
 \mathcal I_y(\xi)=P_y\xi P_y,\qquad
 p_y=\Tr(\xi P_y),\qquad
 \xi_y=\mathcal I_y(\xi)/p_y\quad(p_y>0).
 \label{eq:ii-tins-instrument}
\end{equation}
Its nonselective channel preserves trace and leaves $\Sigma$ invariant.
\end{proposition}
\begin{proof}
Cross-products of distinct $P_y$ vanish. Thus
$\mathbb U^*\mathbb U=\sum_yP_y\otimes R_y^*R_y=I$, and likewise
for the other product. Pointer orthonormality gives $V^*V=\sum_yP_y=I$;
applying this to both sides of $\widehat\xi$ proves recovery. Compressing
the pointer to $|y\rangle$ gives \eqref{eq:ii-tins-instrument}.
Moreover, $\sum_y\Tr(P_y\xi P_y)=\Tr\xi$. Since every $P_y$ commutes
with $\Sigma$, we also have $\sum_yP_y\Sigma P_y=\Sigma$.
\end{proof}

For a state with coherences between marks, forgetting the pointer may
change them. The recovery above uses the joint state, which retains them.
Reading an outcome and retaining the complete apparatus are distinct operations.

Let $J|j\rangle=j|j\rangle$ and $H_{\rm sp}=f(J)\otimes I_M$, with
$f:\mathbb N_0\to\mathbb R$. Its self-adjoint domain is
\[
 \mathcal D(H_{\rm sp})=
 \left\{\psi:\sum_{j,a}|f(j)|^2|\psi_{j,a}|^2<\infty\right\}.
\]
The mark projectors reduce this domain and commute with the spectral
projections of $H_{\rm sp}$. On that domain,
\begin{equation}
 (H_{\rm sp}\otimes I_3)V=VH_{\rm sp},\qquad
 V^*(H_{\rm sp}\otimes I_3)V=H_{\rm sp}.
 \label{eq:ii-tins-energy}
\end{equation}
Indeed, applying $H_{\rm sp}$ to each summand of
\eqref{eq:ii-tins-isometry} allows it to be commuted with $P_y$; the
squared norm of the sum equals the sum of the squared sector norms.
The second identity follows from $V^*V=I$. Energy expectations are
therefore preserved whenever they exist, including the realization
$f(j)=\epsilon_a(j+1/2)$. A pointer Hamiltonian requires separate
specification: if its three levels are degenerate, the coupling also
preserves total energy. This identity assigns no zero cost to physically
preparing or resetting the apparatus.
% END THERMAL_R03_02

% BEGIN THERMAL_R03_03
\subsection{Tritic transport with remainder and carry}
\label{subsec:ii-tins-tritico}

Division $j=3w+b$, $0\le b<3$, defines the unitary
$U_3|j\rangle=|w,b\rangle$. Let $W=U_3\otimes I_M$. The transported marks are
\begin{align}
 P_A'&=\sum_{n=0}^2|7+n,2\rangle\langle7+n,2|\otimes Q_n,\nonumber\\
 P_B'&=|0,0\rangle\langle0,0|\otimes|\chi_B\rangle\langle \chi_B|,
 \qquad P_C'=I-P_A'-P_B'.
 \label{eq:ii-tins-tritic-marks}
\end{align}

\begin{proposition}[Naturality of the apparatus]
\label{prop:ii-tins-tritic}
If $V'$ is the apparatus for the transported marks, then
\begin{equation}
 V'W=(W\otimes I_3)V,\qquad
 \mathcal I_y'(W\xi W^*)=W\mathcal I_y(\xi)W^*.
 \label{eq:ii-tins-square}
\end{equation}
Probabilities and relative states are transported without loss.
\end{proposition}
\begin{proof}
Euclidean division gives $P_y'W=WP_y$ for each $y$. Substituting these
three identities in the sum defining $V'$ proves the first square;
multiplication by the adjoints proves the second. Unitary invariance of
the trace proves the probability statement, and division by the same
positive probability proves the statement for relative states.
\end{proof}

The transported density retains quotient and remainder separately:
\begin{equation}
 W\Sigma W^*=\rho_{r^3}\otimes\sigma_r\otimes I_M/M,
 \qquad \sigma_r=\frac{\operatorname{diag}(1,r,r^2)}{1+r+r^2}.
 \label{eq:ii-tins-tritic-state}
\end{equation}
Indeed, $(1-r)r^{3w+b}=(1-r^3)(r^3)^w r^b/(1+r+r^2)$.
For the unilateral shift $S|j\rangle=|j+1\rangle$,
\begin{equation}
 U_3SU_3^*=I\otimes\bigl(|1\rangle\langle0|+|2\rangle\langle1|\bigr)
 +S_w\otimes|0\rangle\langle2|.
 \label{eq:ii-tins-carry}
\end{equation}
The cases $b=0,1$ increment the remainder; $b=2$ resets the remainder to
zero and increments $w$. This proof retains the carry. The shift may
change the mark: \eqref{eq:ii-tins-square} expresses independence of the
order of observation and reversible representation change, not immobility
under the dynamics.
% END THERMAL_R03_03

% BEGIN THERMAL_R03_04
\subsection{Equivalent embeddings and memory enlargement}
\label{subsec:ii-tins-inclusiones}

Let $\widetilde Q_n$ be other projectors of the same ranks and
$\widetilde \chi_B$ another unit vector. At each degree $j_n$, a unitary
$R_{j_n}$ takes $\operatorname{ran}Q_n$ to
$\operatorname{ran}\widetilde Q_n$: send orthonormal bases of those
subspaces and their complements to one another. Choose likewise
$R_0\chi_B=\widetilde \chi_B$ and complete with unitaries at the remaining
degrees. Then $R=\bigoplus_{j\ge0}R_j$ is unitary, preserves $\Sigma$,
and commutes with $H_{\rm sp}$. The identities
\begin{equation}
 \widetilde P_yR=RP_y,\qquad
 \widetilde V R=(R\otimes I_3)V
 \label{eq:ii-tins-inclusion}
\end{equation}
hold block by block. Conjugation and cyclicity of the trace show that
they preserve probabilities, spectral energy, and $b_*$. A single
degree-independent memory unitary need not realize every embedding
simultaneously. Any additional operator is transported by the same $R$;
this conjugation does not assert that its matrix remains unchanged.

For an enlargement, let $L:\mathbb C^M\to\mathbb C^{M'}$ be an isometry
and let $J_L=I\otimes L$. Define
$P_A^+=J_LP_AJ_L^*$, $P_B^+=J_LP_BJ_L^*$, and
$P_C^+=I-P_A^+-P_B^+$. Since $J_L^*J_L=I$, all three marks satisfy
$P_y^+J_L=J_LP_y$. Consequently,
\begin{equation}
 V^+J_L=(J_L\otimes I_3)V,
 \qquad \Sigma_L=J_L\Sigma J_L^*=\rho_r\otimes LL^*/M.
 \label{eq:ii-tins-extension}
\end{equation}
Isometric transport preserves the original probabilities. If the uniform
density of the whole enlarged ambient space is also prepared,
$\Sigma^+=\rho_r\otimes I_{M'}/M'$, the preparation changes and
\begin{equation}
 p_A^+=\frac{M}{M'}p_A,\qquad
 p_B^+=\frac{M}{M'}p_B,\qquad p_A^+/p_B^+=b_*.
 \label{eq:ii-tins-uniform-extension}
\end{equation}
The proof again takes the trace of each degree's scalar over the same
ranks. On $A$ and $B$, normalization cancels $M/M'$ and gives the
isometric images of the original relative states; the complement receives
the added weight. Preservation of the apparatus, equality of the intensity,
and equality of the complete state are thereby distinguished.
% END THERMAL_R03_04

% BEGIN THERMAL_R03_05
\subsection{Complete preparation, entropy, and conditioning}
\label{subsec:ii-tins-condicionamiento}

Let $\rho$ be a state of a finite trial system, $H=H^*$ its inherited
Hamiltonian, and $\rho_0$ a fixed pure preparation. With $P=P_A$, define
\begin{equation}
 \Omega_\rho=(P\Sigma P)\otimes\rho
 +((I-P)\Sigma(I-P))\otimes\rho_0,
 \qquad \Omega_0=\Sigma\otimes\rho_0.
 \label{eq:ii-tins-preparation}
\end{equation}
This preparation retains the entire spectral complement. It differs from
the auxiliary product in \eqref{eq:ii-ter27-flag}, although it realizes
the same two functionals on the trial system.

\begin{theorem}[Information per reference and conditional energy]
\label{thm:ii-tins-readings}
The trace of $\Omega_\rho$ is one, and its partial trace over the trial
system recovers $\Sigma$. Moreover,
\begin{align}
 \mathcal S_B(\rho)&:=
 \frac{s(\Omega_\rho)-s(\Omega_0)}{p_B}=b_*s(\rho),
 \label{eq:ii-tins-information}\\
 \rho_A&:=\frac{\Tr_{\mathcal K}[(P\otimes I)\Omega_\rho(P\otimes I)]}{p_A}
 =\rho,\qquad
 \mathcal E_A(\rho)=\Tr(\rho_AH)=\Tr(\rho H).
 \label{eq:ii-tins-conditioned}
\end{align}
\end{theorem}
\begin{proof}
The positive summands have traces $p_A$ and $1-p_A$. Since
$[P,\Sigma]=0$, partial trace over the trial system gives
$P\Sigma P+(I-P)\Sigma(I-P)=\Sigma$. Let $\lambda_i$ be the eigenvalues
of the marked block and $p_j$ those of $\rho$. Their products replace
$\lambda_i$; the entropy increment of that block is
\[
 -\sum_{i,j}\lambda_ip_j\ln(\lambda_ip_j)
 +\sum_i\lambda_i\ln\lambda_i=p_As(\rho).
\]
The complement is unchanged. The sums are legitimate because
$s(\Sigma)=-\ln(1-r)-r\ln r/(1-r)+\ln M<\infty$ and the trial
system is finite. Division by $p_B$ proves
\eqref{eq:ii-tins-information}. Compression to $P\otimes I$ and partial
trace give $p_A\rho$; division by $p_A>0$ proves
\eqref{eq:ii-tins-conditioned} and its energy expectation.
\end{proof}

Conditioning preserves the observational operation of the corpus. If
$U_a$ couples the state to the initial record $m_0$, the output and its
fibre are
\[
 \operatorname{Out}_a(x)=q_a(\operatorname{pr}_M U_a(x,m_0)),\qquad
 \Omega_{a,y}=\{x:\operatorname{Out}_a(x)=y\}.
\]
For $\mu_a(\Omega_{a,y})>0$, the measure is restricted to that fibre:
\begin{equation}
 \mu_{a,y}(D)=\frac{\mu_a(D\cap\Omega_{a,y})}{\mu_a(\Omega_{a,y})}.
 \label{eq:ii-tins-fiber}
\end{equation}
When these events belong to a common projection-valued measure and
$\mu_a(D)=\Tr(\omega P_D)$, we have
$P_DP_y=P_{D\cap\Omega_{a,y}}$. Cyclicity of the trace proves
\[
 \Tr\!\left(\frac{P_y\omega P_y}{\Tr(\omega P_y)}P_D\right)
 =\mu_{a,y}(D).
\]
The equality extends to bounded observables of that algebra by spectral
sums and uniform approximation. Applied to event $A$ in
\eqref{eq:ii-tins-preparation}, it gives exactly the preceding conditional
mean energy. It requires no dynamical ergodicity assumption. Any additional
purification is retained in its own factor; the entropy computed here
belongs to the specified observable algebra.
% END THERMAL_R03_05

% BEGIN THERMAL_R03_06
\subsection{Characterization by additivity and refinement}
\label{subsec:ii-tins-aditividad}

The preceding evaluation admits a characterization fixing its weight
without postulating $b_*$ as a coefficient. For a fixed trial state $\rho$,
consider a contribution $F_\rho(x)$ for an alternative of received weight
$x\in[0,1]$. In the relative reading, this weight is measured against
reference $B$; only alternatives within the stated domain are considered.
Require dependence on the weight rather than the alternative's name;
additivity $F_\rho(x+y)=F_\rho(x)+F_\rho(y)$ for $x+y\le1$;
positivity; and normalization $F_\rho(1)=s(\rho)$. The domain must
realize these partitions or the ordered refinements described below.

\begin{theorem}[Uniqueness of the additive informational weight]
\label{thm:ii-tins-additive}
On the complete scalar domain, the four conditions above imply
$F_\rho(x)=xs(\rho)$. The same conclusion holds for a realization of
alternatives with equiponderant refinements of weights $3^{-n}$ and,
for each mark $P$, approximations $P_n^-\le P\le P_n^+$ whose weights
converge to that of $P$ and whose contributions belong to those refinements.
\end{theorem}
\begin{proof}
Additivity gives $F_\rho(0)=0$. If $x\le y$, then
$F_\rho(y)-F_\rho(x)=F_\rho(y-x)\ge0$, so the reader is monotone.
Partitioning the unit into $3^n$ equal alternatives gives
$F_\rho(3^{-n})=3^{-n}s(\rho)$, and joining $k$ of them gives
$F_\rho(k3^{-n})=k3^{-n}s(\rho)$. For $0\le x<1$, take
\[
 \ell_n=3^{-n}\lfloor3^nx\rfloor,\qquad u_n=\ell_n+3^{-n}.
\]
They bound $x$ within $[0,1]$. Monotonicity implies
$\ell_ns(\rho)\le F_\rho(x)\le u_ns(\rho)$; their difference
$3^{-n}s(\rho)$ tends to zero. This proves the formula, including
$s(\rho)=0$; $x=1$ is the normalization. For ordered marks, repeat the
argument with $P_n^-$ and $P_n^+$, since positive additivity implies
isotonicity. No differentiable continuity of the reader is assumed.
\end{proof}

In particular, the contribution of $A$ is $b_*s(\rho)$ when its relative
measure is retained. The characterization also applies to the normalized
mark of \eqref{eq:ii-ter27-flag}. There, the functional
$-\Tr[(P\otimes I)(\widehat\eta\otimes\rho)(I\otimes\ln\rho)]$
satisfies the conditions by linearity of the trace and assigns $s(\rho)$
to the certain mark.

Refinement requires an effective realization: the rank-one projector
$P_B$ contains no three nonzero equiponderant subprojectors. An auxiliary
enlargement can realize refinements, but must transport the measure and
the reading. Likewise, three descendants of a cylinder do not imply
equal weights merely from their number. For a rational weight, an
equiponderant partition with the corresponding denominator suffices if
it exists in the realization; the ternary proof does not assume that
particular denominator.

The assumptions isolate precisely the freedoms excluded. The reader
$f(x)=x^2$ respects products but gives
$f(x+y)-f(x)-f(y)=2xy$; for $x=1/3$, $y=1/5$, the defect is $2/15$.
Without normalization, $f(x)=cx$ preserves positivity and additivity.
Without identifying the received measure, $(1/3,2/3)$ and $(1/2,1/2)$
are distinct positive, additive, normalized measures on two alternatives.
The fixed entropy of the label record remains separate: the theorem
does not remove the mixing term from a joint entropy.
% END THERMAL_R03_06

% BEGIN THERMAL_R03_07
\subsection{Realization by visits between returns}
\label{subsec:ii-tins-retornos}

Relative intensity also has a precise dynamical interpretation. Let
$(X,\mu,T)$ be an invertible, ergodic, measure-preserving probability
system with events $A,B$ of measures $p_A,p_B>0$. For $x\in B$, let
\[
 R_B(x)=\min\{k\ge1:T^kx\in B\},\qquad
 C_A(x)=\sum_{k=0}^{R_B(x)-1}\mathbf1_A(T^kx).
\]
The excursion retains its complete sequence; $C_A$ is only one observable
of that sequence.

\begin{theorem}[Reward per return to a reference]
\label{thm:ii-tins-return}
With $\mu_B=\mu|_B/p_B$, for every $f\in L^1(\mu)$,
\begin{equation}
 \mathbb E_{\mu_B}\!\left[\sum_{k=0}^{R_B-1}f\circ T^k\right]
 =\frac1{p_B}\int_Xf\,d\mu.
 \label{eq:ii-tins-return}
\end{equation}
In particular, $\mathbb E_{\mu_B}C_A=p_A/p_B=b_*$.
\end{theorem}
\begin{proof}
Let $B_k=\{x\in B:R_B(x)>k\}$. Recurrence, invertibility, and
ergodicity ensure that the levels $T^kB_k$, $k\ge0$, partition $X$
up to a null set: every point has a most recent visit to $B$ in its past.
For $f\ge0$, measure preservation and monotone convergence give
\[
 \int_B\sum_{k=0}^{R_B-1}f(T^kx)\,d\mu
 =\sum_{k\ge0}\int_{T^kB_k}f\,d\mu=\int_Xf\,d\mu.
\]
For $f\in L^1$, applying the identity first to $|f|$ establishes
integrability; the positive and negative parts can then be subtracted.
Division by $p_B$ proves the formula; $f=\mathbf1_A$ gives its consequence.
\end{proof}

For a finite irreducible chain with matrix $P$ and stationary distribution
$\pi>0$, there is also an algebraic proof. Starting at the reference
state $B$, let $v_i$ be the expected number of visits to $i$ before the
first positive return. That return has finite expectation; summing
transitions up to it shows that arrivals and departures have equal
expectations, hence $vP=v$. Since $v_B=1$, uniqueness of the stationary
distribution gives $v_i=\pi_i/\pi_B$. Summing over $A$ reproduces
\eqref{eq:ii-tins-return}. Grouping the spectral distribution into $B$,
the three sectors of $A$, and the complement provides a finite realization
with the exact ratio \eqref{eq:ii-tins-relative}.

If every visit to $A$ carries a trial $\rho$, independent of the others
conditional on the excursion record, the trials form
$\rho^{\otimes C_A}$. Tensor additivity gives
\begin{align}
 \mathbb E_{\mu_B}s(\rho^{\otimes C_A})&=b_*s(\rho),\nonumber\\
 \frac{\mathbb E_{\mu_B}[C_A\Tr(\rho H)]}
      {\mathbb E_{\mu_B}C_A}&=\Tr(\rho H).
 \label{eq:ii-tins-visit-readings}
\end{align}
The second expression is the ratio of total energy to total visits,
not an average of ratios per excursion, which would be undefined when
$C_A=0$. Accumulated energy per return equals $b_*\Tr(\rho H)$.

Information per return and energy per visit thus realize the two
denominators of \eqref{eq:ii-tins-information}--
\eqref{eq:ii-tins-conditioned}. Correlated trials require their joint
conditional entropy. A noninvertible process admits this formulation
on its stationary bilateral extension; without ergodicity, the domain
must be restricted to the support swept out by the return section.
These dynamical assumptions are unnecessary for the operator theorem
\ref{thm:ii-tins-readings}. Neither a stationary test chain is identified
by itself with the TPK update, nor $R_B$ with the material period $108t_0$.
% END THERMAL_R03_07

% BEGIN THERMAL_R03_08
\subsection{Common reference and correlations between systems}
\label{subsec:ii-tins-comun}

One reference and $m$ independent systems produce the record
$\widehat\eta\otimes\rho_1\otimes\cdots\otimes\rho_m$, with a
single mark $P\otimes I$. Its weight remains $b_*$. For
$H_{\rm tot}=\sum_jH_j$, with the appropriate tensor extensions,
\begin{equation}
 \mathcal S_{P,\rm tot}=b_*\sum_js(\rho_j),\qquad
 \mathcal E_{P,\rm tot}=\sum_j\Tr(\rho_jH_j).
 \label{eq:ii-tins-common}
\end{equation}
The proof uses trace factorization and
$\ln(\rho_1\otimes\rho_2)=\ln\rho_1\otimes I+I\otimes\ln\rho_2$
on the support, iterated over the factors. For a correlated state
$\rho_{1\ldots m}$, the same mark gives $b_*s(\rho_{1\ldots m})$.
For two systems,
\begin{equation}
 \mathcal S_{P,1}+\mathcal S_{P,2}-\mathcal S_{P,12}
 =b_*I(1:2)_\rho,
 \label{eq:ii-tins-mutual}
\end{equation}
by the definition of mutual information. An interacting Hamiltonian also
retains a single reference even when its Gibbs state does not factorize.

Two independent references with joint mark $P_1\otimes P_2$ have weight
$b_1b_2$; for correlated references the weight is
$\Tr[\omega_{R_1R_2}(P_1\otimes P_2)]$ and need not factorize.
Copying a label does not create independence. Indeed, the isometry
\[
 V_Pv=Pv\otimes|1\rangle+(I-P)v\otimes|0\rangle
\]
satisfies $V_P^*(P\otimes|1\rangle\langle1|)V_P=P$; the trace in
the transported state preserves $b_*$, not $b_*^2$. More generally,
for an isometry $L$,
$\Tr[(L\omega L^*)(LPL^*)]=\Tr(\omega P)$ because $L^*L=I$.
Returning while retaining the same record does not automatically raise
its weight to a power. A new survival selection uses its own mark.
% END THERMAL_R03_08

% BEGIN THERMAL_R03_09
\subsection{Thermometric selection and composition with action and horizon}
\label{subsec:ii-tins-seleccion}

Fix the reference and a finite nonscalar Hamiltonian $H$. For
$\rho_\beta=e^{-\beta H}/Z(\beta)$, $\beta>0$, the additive theorem
and conditional energy produce the pair
$\mathcal S_B=b_*s(\rho_\beta)$, $\mathcal E_A=\Tr(\rho_\beta H)$.
Writing $E=\Tr(\rho_\beta H)$, the finite spectral sum gives
$E'=-\operatorname{Var}_{\rho_\beta}(H)<0$ and $(\ln Z)'=-E$.
Differentiating $s=\beta E+\ln Z$ gives $s'=\beta E'$, hence
\begin{equation}
 \frac{d\mathcal S_B}{d\mathcal E_A}=b_*\beta,\qquad
 \Theta_{B,A}=\frac1{b_*\beta},\qquad B_*(u_\Theta)=b_*u_E.
 \label{eq:ii-tins-temperature}
\end{equation}
The positive bases satisfy $u_S=u_E/u_\Theta$. The last equation
determines the unique positive linear map sending
$\Theta_{B,A}u_\Theta$ to $\beta^{-1}u_E$. The balance
\[
 \Tr(\rho H)-b_*\Theta s(\rho)-
 \bigl[\Tr(\rho_\Theta^*H)-b_*\Theta s(\rho_\Theta^*)\bigr]
 =b_*\Theta D(\rho\Vert\rho_\Theta^*)
\]
with $\rho_\Theta^*\propto e^{-H/(b_*\Theta)}$ retains the unique
minimum proved in \ref{cor:ii-ter27-variational}. Additive selection
now characterizes the weight entering that balance.

With $\epsilon_a=h_a/(108t_0)=\hbar_ac/L_\alpha$,
$q_A=\gamma a_0L_\alpha^2$, and
$\mathcal T_{I/A}=\epsilon_a/q_A$, the thermal sections are
\begin{align}
 b_*\Theta_{{\rm clk},P,a}
 &=\frac{h_a}{108t_0\ln3}
 =\frac{\mathcal T_{I/A}q_A}{\ln3},\nonumber\\
 b_*\Theta_{H,P}(R)&=\frac{\hbar_ac}{2\pi R},&
 \frac{\Theta_{H,P}(R)}{\Theta_{{\rm clk},P,a}}
 &=\frac{L_\alpha}{R}\frac{\ln3}{2\pi}.
 \label{eq:ii-tins-action-area}
\end{align}
These are coordinate equations for $B_*$ in the transported bases.
The first uses the section $\beta_{{\rm clk},a}=\ln3/\epsilon_a$;
the second uses the scale $\tau_R=\hbar_ac/(2\pi R)$. Their quotient
proves the temperature ratio. The radial relation
$G_a\hbar_a=c^3L_\alpha^2$ implies
\begin{equation}
 G_ab_*\Theta_{{\rm clk},P,a}\ln3=c^4L_\alpha.
 \label{eq:ii-tins-gravity}
\end{equation}
At the horizon, $\mathcal S_B=b_*s$ and
$\Theta_{H,P}=\tau_R/b_*$ convert the previous balance into
\[
 \delta E_R-\Theta_{H,P}\delta\mathcal S_B
 =\tau_RD(\sigma\Vert\rho_A).
\]
The composition retains Barbero's area resolution, sector incidence,
action, and the energy origin: cancellation of scales preserves the same
energy rather than adding a metrological parameter to the constructor.

The selection has a specified domain. Conditioning both channels gives
$(E,s)$; counting both per reference gives $(b_*E,b_*s)$. In both cases
the derivative is $\beta$. The pair characterized here is $(E,b_*s)$,
using energy per visit and information per reference. Preservation of
the record allows all three observations but does not identify them.
Nor do energy products alone select the weight: for $p>0$, replacing
$b_*$ by $b_*^p$ and dividing temperatures by the same relative factor
preserves those products and serial composition; the quadratic reader,
however, violates the additivity proved above.

The mark instrument, its recovery and transport, the additive informational
contribution, conditional energy, and this realization's transducer are
thus determined. Identifying it with a thermometric instrument of the
complete constructor requires transporting its events, measure, and pair
of normalizations to $A$ and $B$ while retaining refinement and memory.
Existence of the apparatus does not replace that particular identification.
If the reference varies with $\beta$, its derivatives also enter; if the
projection does not commute with $\Sigma$, preparation
\eqref{eq:ii-tins-preparation} preserves its block diagonal rather than
necessarily the original density. These conditions specify the scope
of the identities without altering the preceding results on action,
area, capacity, and transport.
% END THERMAL_R03_09
